%! Author = chtompki
%! Date = 1/14/26
\documentclass[12pt]{amsart}
% Default margins are too wide all the way around. I reset them here
\setlength{\hoffset}{-.75in}
\setlength{\textwidth}{6.5in}
\setlength{\voffset}{-.5in}
\setlength{\textheight}{9.0in}
\usepackage{amsmath}
\usepackage{leftindex}
\usepackage{amssymb}
\usepackage{amsthm}
\usepackage{enumitem}
\usepackage{setspace}

\newenvironment{solution}
{\begin{proof}[Solution]}
{\end{proof}}

\newtheorem{theorem}{Theorem}

\title{Combinatorics\\Catalan Numbers Homework}
\author{Rob Tompkins}
\date{\today}

\begin{document}
    \maketitle
    \date{\today}
    \begin{onehalfspace}
        \begin{theorem}
            Let $C_n$ represent the Catalan numbers, with
            \[
                C_n = \frac{1}{n+1}\binom{2n}{n}.
            \].
            We define the Catalan numbers as the number of valid ways to arrange $n$-pairs of left-right
            parentheses where every opening bracket has a subsequent closing bracket, and they are nested correctly.
            We want to show that $C_n$ as defined above, also represents the number of binary trees with $n$ vertices.
        \end{theorem}

        \begin{proof}
            Our proof technique is the mechanics of reduction.
            We reduce one problem to another based on definition.
            Our goal will be to show that there is a bijective correspondence between the two definitions and
            their outcomes.

            The fundamental idea behind the proof will be that every binary tree is determined by a root and a pair
            of smaller trees.
            So counting binary trees in this fashion is a split of counting two independent sub-Catalan objects.

            Supose $\Phi$ defines a map from binary trees to balanced parentheses strings.
            We do the following iteratively: (a) the empty tree maps to the empty string, and (b) if a tree has a root,
            a left subtree $L$, and a right subtree $R$, then we write $\Phi(T) = (\Phi(L))\Phi(R)$.
        \end{proof}
    \end{onehalfspace}
\end{document}
